%% kapa: exam=1 round=28 year=2017 period=2 form=A subject=appraisal_law subject_ko=감정평가관계법규 seq=33 number=33
%% source: Q-Net 감정평가사 기출문제 (cst00309) artlSeq=5199603 file=extracted/1008/5199603/2017년 제28회 감정평가사 1차시험 문제지/2교시 A형.pdf
%% conversion: kordoc 4.24.51 -> tools/gen1.py (auto); review=proofread-sonnet; flags=none
\begin{question}{33}
부동산등기법령상 등기부 등에 관한 설명으로 옳지 않은 것은?
\begin{choices}
  \item 등기부는 영구히 보존하여야 한다.
  \item 법관이 발부한 영장에 의하여 신청서나 그 밖의 부속서류를 압수하는 경우에는 이를 등기소 밖으로 옮길 수 있다.
  \item 1동의 건물을 구분한 건물에 있어서는 1동의 건물에 속하는 전부에 대하여 1개의 등기기록을 사용한다.
  \item 등기기록의 부속서류에 대하여는 이해관계 있는 부분만 열람을 청구할 수 있다.
  \item 등기관의 중복등기기록 정리는 실체의 권리관계에 영향을 미친다.
\end{choices}
\end{question}
\answer{5}
